\documentclass[10pt,a4paper]{amsart}
\usepackage[margin=19mm]{geometry}
\usepackage{amsmath,amssymb,mathtools}
\usepackage[scr=rsfs,cal=euler]{mathalfa}
\usepackage{xfrac}
\usepackage[unicode,hidelinks,bookmarks=false]{hyperref}
\newcommand{\ie}{i.e.\ }
\newcommand{\cf}{cf.\ }
\newcommand{\resp}{respectively\ }
\newcommand{\Subsection}{\subsection}
\providecommand{\nobreakdash}{\nobreak\hbox{-}}
\setlength{\emergencystretch}{2em}
\multlinegap0pt
\usepackage{amssymb}
\newtheorem{thm}{Theorem}
\newtheorem{lemma}{Lemma}
\def\R{\mathbb R}
\newcommand{\arcs}[1]{\stackrel{\rotatebox{90}{\negthinspace\hspace{1pt}\big)}}{#1}} 
\def\d{\delta}
\def\om{\omega}
\title{Lipschitzness of the width and diameter functions of convex bodies in $\R^n$}
\author{Oleg Mushkarov, Nikolai Nikolov,, Pascal J. Thomas}
\date{Source: arXiv:2406.12537v1 (CC BY 4.0)}
\begin{document}
\maketitle

\noindent\emph{Source and licence.} Lipschitzness of the width and diameter functions of convex bodies in $\R^n$. Version arXiv:2406.12537v1 (CC BY 4.0), distributed by arXiv under the Creative Commons Attribution 4.0 International licence, http://creativecommons.org/licenses/by/4.0/, as recorded in the arXiv licence metadata for this exact version and shown on the abstract page https://arxiv.org/abs/2406.12537v1. Full licence notice: \texttt{licenses/arxiv-2406.12537-CC-BY-4.0.txt}.\par\smallskip

\noindent\textit{Editorial scope.} Counted unit: the article's thm main (source0.tex lines 147--154). The article's complete original proof of it is delivered with it (source0.tex lines 156--294, one \texttt{proof} environment of 139 lines). No mathematics is written, repaired or re-derived: the statement and the proof are the article's own lines, in the article's own notation, and no retained line is substituted or re-flowed.  No further source-local context is needed: every cross-reference of every retained line resolves inside the two delivered spans.  Nothing was omitted from any retained span, so the package carries no omission record.  No retained line contains a citation, so the wrapper carries no bibliography and no bibliography entry.  No retained line contains a figure environment, an image-inclusion command or drawing code, so no image is delivered and the package declares no asset dependency.  Editorial formatting only, in addition: an AMS \texttt{amsart} preamble that restates the article's own declarations and macros used by the retained lines, namely the environments \texttt{thm}, \texttt{lemma} on separate counters, together with the standard packages those lines need.  The retained environments print in this document as Theorem 1, Lemma 1 in order of appearance, whereas the article prints the same items with the numbering of its own full document.  The complete native source of the article as distributed in the arXiv e-print for this exact version is bundled unchanged as \texttt{sources/160956/source0.tex}.
\begin{thm}
\label{main} For any convex body $K$ in $\R^n$ and $u, v\in
S^{n-1}$:
\begin{enumerate}
\item $\Delta w_K(u,v)\le M_K$;
\item $\Delta d_K(u,v)\le N_K$.
\end{enumerate}
\end{thm}

\begin{proof}

Since $K$ is the intersection of convex polytopes $K_1\supset K_2\supset\dots$
and $w_{K_i}\downarrow w_K,$ $d_{K_i}\downarrow d_K,$ $\d_{K_i}\downarrow\d_K,$ and
$\om_{K_i}\downarrow\om_K,$ we may assume that $K$ is a convex polytope.
\smallskip

{\it Proof of Part (1).} Since $w_K$ is already known to be
continuous, it will be enough to prove the inequality outside of a
set of dimension $n-2$ on the unit sphere. Let $E$ be an edge of
$K$ (non-trivial line segment contained in $\partial K$ such that
its relative interior is not contained in the relative interior of
any face of larger dimension), and $H_E$ the  hyperplane of
vectors orthogonal to $E$. For further reference, call those
\emph{exceptional hyperplanes}. Let $u\in \mathcal W_K := S^{n-1}
\setminus \bigcup_{E \mbox{ \small edge of }K} H_E$.  Then $u$ is
not orthogonal to any face of $\partial K$ of any dimension $d$,
$1\le d\le n-1$. For such a vector $u$, each of the two supporting
hyperplanes $H_1$ and $H_2$ orthogonal to $u$ intersects $K$ at a
single point, denoted respectively by $A$ and $B$.

\begin{lemma}
\label{limw}
For any $u \in \mathcal W_K$ and $A, B$ chosen as above,
\[
w_K'(u):= \limsup_{u',u''\to u} \Delta w_K(u',u'')
= \sqrt{|AB|^2-w^2_K(u)} \le M_K.
\]
\end{lemma}
\begin{proof}
For any $v$ close enough to $u$, the corresponding supporting
hyperplanes intersect $K$ at the same points $A$ and $B$, so that
$w_K(v)= |\langle AB, v\rangle|$ and
\[
\frac{|w_K(u')-w_K(u'')|}{\|u'-u''\|} = \left| \left\langle AB, \frac{u'-u''}{\|u'-u''\|}\right\rangle\right|.
\]
Note that $\langle u'-u'', (u'+u'')/2 \rangle=0$ and that $(u'+u'')/2$ tends to $u$.
Passing to a subsequence, we may assume that $\frac{u'-u''}{\|u'-u''\|} \to \tilde u$, with
$\tilde u \in S^{n-1}$, $\langle \tilde u, u \rangle=0$.
Let $\tilde A$ be the
unique point in  the supporting hyperplane going through $B$
such that $A\tilde A$ is parallel to $u$; then $\langle AB, \tilde u\rangle = \langle \tilde AB, \tilde u\rangle$.
So, observing that $\rho(u',u'') \sim \|u'-u''\|$ as $\|u'-u''\|\to0$,
\begin{multline*}
\lim_{u',u''\to u} \Delta w_K(u',u'') =
\lim_{u',u''\to u}\left| \left\langle AB, \frac{u'-u''}{\|u'-u''\|}\right\rangle\right|
\\
= \langle \tilde AB, \tilde u\rangle \le |\tilde AB| = \sqrt{|AB|^2-|A\tilde A|^2} =
\sqrt{|AB|^2-w^2_K(u)}.
\end{multline*}
We can obtain the case of equality above if we choose $u'=u$ and $u''$ so that $\tilde u$
is parallel to $(\tilde A B)$, which shows that the limes superior is equal to the righthand
side.
\end{proof}

Any two vectors $u,v$ in the same connected component of
$\mathcal W_K$ can be joined
within the component by a geodesic (arc of a great circle). If
the inequality in Theorem \ref{main} (1) was violated, by dichotomy, we could find
arbitrarily small arcs $\arcs{u'u''}$
within the geodesic segment where the inequality
would be violated, but since $\rho(u',u'')  \sim \|u'-u''\|$ as both quantities tend
to $0$, this would violate Lemma \ref{limw}.

If $u$ or $v$ lies on an exceptional hyperplane, and the other in
an adjacent component, we obtain the same result by continuity. If
$u$ and $v$ belong to different components, the shorter of the two
arcs on the great circle containing them will intersect the
exceptional hyperplanes in a finite number of points, and we
conclude again by continuity.  If $u$ and $v$ are on distinct
exceptional hyperplanes, the same method works.  Finally, if $u$
and $v$ lie on the same exceptional hyperplane, we may perturb one
of them slightly, apply the previous estimate, and conclude again
by continuity.

\smallskip

{\it Proof of Part (2).} As in Part (1), we exclude a finite set
of exceptional vector hyperplanes from the unit sphere. This time
those are the hyperplanes $H_D$ generated by a set $D$ of vertices
of $K$ with $\# D=n$. When $u \in \mathcal R_K := S^{n-1}
\setminus \bigcup_{D\subset V(K), \# D=n} H_D $, we can choose a
diametral chord for $u$ passing through a vertex $A$ of $K$; its
other extremity $B$ will have to be in the interior of an
$(n-1)$-dimensional face of $K$. Denote by $\tilde d_K(u)$ the
distance between the parallel supporting hyperplanes of $K$, $H_1$
and $H_2$, that contain $A$ and $B$ respectively.

\begin{lemma}
\label{limd}
For any $u\in \mathcal R_K$,
\[
d_K'(u):= \limsup_{u',u''\to u} \Delta d_K(u',u'')
=
\frac{d_K(u)}{\tilde d_K(u)}\sqrt{d^2_K(u)-\tilde d^2_K(u)}\le N_K.
\]
\end{lemma}

\begin{proof}
There is a neighborhood of $u$, $\mathcal N\subset  \mathcal R_K$,
such that for any $t\in \mathcal N$, a diametral chord for $t$
also passes through $A$ and the face containing $B$. Assume $u',
u''\in \mathcal N$.

Let $H$ be the \emph{strip} of $\mathbb R^n$ determined by $H_1$
and $H_2$ (that is, the intersection between two distinct half
spaces determined by $H_1$ or $H_2$ that contains $K$). Consider
the points $B' \in H_2$ (resp. $B''\in H_2$) such that $AB'$(resp.
$AB''$) is parallel to $u'$ (resp. $u''$), then $d_K(u')=|AB'|$,
$d_K(u'')=|AB''|$. Assume as we may that $|AB''|\ge |AB'|$. In
$\triangle AB'B''$, let $h$ be the distance from $A$ to the line
$(B'B'')$ and  $\alpha, \beta', \beta''$ be the angles at the
vertices $A, B', B''$ respectively. Again we use that
$\rho(u',u'') \sim \|u'-u''\|$. Then by the law of sines and
$\alpha+ \beta'+ \beta''=\pi$,
\begin{multline*}
\frac{|AB''|- |AB'|}{\|u'-u''\|} = \frac{|AB'|(\sin \beta' - \sin \beta'')}
{ (\sin \beta'')( 2 \sin \alpha/2) } = |AB'| \frac{\sin \frac{\beta'-\beta''}
2  \cos \frac{\beta'+\beta''}2}{ \sin \beta''  \sin \alpha/2 }
\\
= |AB'| \frac{\cos (\beta''+  \frac{\alpha}2)}{ \sin \beta'' } \le |AB'| \cot \beta'' =
|AB'| \sqrt{\frac{|AB''|^2}{h^2} -1 } \le  |AB'| \sqrt{\frac{|AB''|^2}{\tilde d_K(u)^2} -1 },
\end{multline*}
since $h\ge \tilde d_K(u)$. Passing to the limit, we obtain that $N_K$
is an upper bound for the limes superior we are considering.

Observe that as $u', u''\to u$, then $\alpha \to 0$, $|AB'|,|AB''|\to d_K(A)$,
but when the unit vectors parallel to $(B'B'')$ admit a limit $v$,
$h\to \mbox{dist}(A, L)$
where $L$ is the line through $B$ parallel to $v$.
If we choose $A' \in H_B$ such that $|AA'|= \tilde d_K(u)$, and $u',u''$
so that $B', B''$ tend to $B$ along the line $(BA')$, then
$h\to  \tilde d_K(u)$, and $$\lim_{u',u''\to u}\Delta d_K(u',u'')=
\frac{d_K(u)}{\tilde d_K(u)}\sqrt{d^2_K(u)-\tilde d^2_K(u)}.$$
\end{proof}

We finish the proof of Part (2) in the same way as that of Part
(1), substituting Lemma \ref{limd} for Lemma \ref{limw}.
\end{proof}

\end{document}
